% Part: normal-modal-logic
% Chapter: completeness
% Section: complete-consistent-sets

\documentclass[../../../include/open-logic-section]{subfiles}

\begin{document}

\olfileid{nml}{com}{ccs}

\olsection{संपूर्ण $\Sigma$-सुसंगत संच}

$\Sigma$ हा मोडल !!{formula}s चा संच आहे असे
समजा---सामान्य मोडल तर्कशास्त्राची स्वयंसिद्धके
किंवा परिभाषक तत्त्वे म्हणून त्यांच्याकडे पाहा.
संच~$\Gamma$ हा $\Sigma$-सुसंगत असेल तेव्हा
आणि तेव्हाच $\Gamma \Proves/[\Sigma] \lfalse$;
म्हणजे प्रत्येक $!A_i \in \Gamma$ असताना
$\Sigma$ पासून~$!A_1 \lif (!A_2 \lif \cdots (!A_n \lif
\lfalse)\dots)$ ची !!{derivation} नसते.
आपण ``कॅनॉनिकल'' प्रतिमान रचू; त्यात प्रत्येक
जग विशिष्ट प्रकारचा $\Sigma$-सुसंगत संच
असेल. तो केवळ $\Sigma$-सुसंगत नसून
महत्तम अर्थाने तसा असेल: प्रत्येक मोडल
!!{formula} चे सत्यत्व तो ठरवेल.
प्रत्येक $!A$ साठी $!A
\in \Gamma$ किंवा $\lnot !A \in \Gamma$:

\begin{defn}
  संच~$\Gamma$ हा \emph{संपूर्ण $\Sigma$-सुसंगत}
  असेल तेव्हा आणि तेव्हाच तो $\Sigma$-सुसंगत
  असेल आणि प्रत्येक~$!A$ साठी
  $!A \in \Gamma$ किंवा $\lnot !A \in \Gamma$
  यांपैकी एक विधान खरे असेल.
\end{defn}

संपूर्ण $\Sigma$-सुसंगत संच~$\Gamma$ चे अनेक
उपयुक्त गुणधर्म आहेत. प्रथम, ते निगामीदृष्ट्या
बंद असतात: $\Gamma
\Proves[\Sigma] !A$ असेल, तर $!A \in \Gamma$.
विशेषतः, !!a{formula}~$!A \in \Sigma$ चा प्रत्येक
दृष्टांतही $\in
\Gamma$ असतो. शिवाय $\Gamma$ मधील सदस्यत्व
विधानीय संयोगकांच्या सत्यत्व-अटींशी जुळते.
आपण ``कॅनॉनिकल प्रतिमान'' परिभाषित करू तेव्हा
हे महत्त्वाचे ठरेल.

\begin{prop}\ollabel{prop:ccs-properties}
  $\Gamma$ हा संपूर्ण $\Sigma$-सुसंगत आहे असे
  समजा. मग:
  \begin{enumerate}
  \item \ollabel{prop:ccs-closed}%
    $\Gamma$ हा~$\Sigma$ च्या सापेक्ष
    निगामीदृष्ट्या बंद आहे.
  \item \ollabel{prop:ccs-sigma}%
    $\Sigma \subseteq \Gamma$.
  \tagitem{prvFalse}{\ollabel{prop:ccs-lfalse}%
    $\lfalse \notin \Gamma$}{}
  \tagitem{prvTrue}{\ollabel{prop:ccs-ltrue}%
    $\ltrue \in \Gamma$}{}
  \item \ollabel{prop:ccs-lnot}%
    $\lnot!A \in \Gamma$ असेल तेव्हा आणि
    तेव्हाच $!A \notin
    \Gamma$.
  \tagitem{prvAnd}{\ollabel{prop:ccs-land}%
    $!A \land !B \in \Gamma$ असेल तेव्हा आणि तेव्हाच
    $!A \in \Gamma$ आणि $!B \in \Gamma$ असतील.}{}
  \tagitem{prvOr}{\ollabel{prop:ccs-lor}%
    $!A \lor !B \in \Gamma$ असेल तेव्हा आणि तेव्हाच
    $!A \in \Gamma$ किंवा $!B \in \Gamma$ असेल.}{}
  \tagitem{prvIf}{\ollabel{prop:ccs-lif}%
    $!A \lif !B \in \Gamma$ असेल तेव्हा आणि तेव्हाच
    $!A \notin \Gamma$ किंवा $!B \in \Gamma$ असेल.}{}
  \tagitem{prvIff}{\ollabel{prop:ccs-liff}%
    $!A \liff !B \in \Gamma$ असेल तेव्हा आणि तेव्हाच
    $!A \in \Gamma$ आणि $!B \in
    \Gamma$ दोन्ही असतील, किंवा
    $!A \notin \Gamma$ आणि $!B \notin \Gamma$
    दोन्ही असतील.}{}
  \end{enumerate}
\end{prop}

\begin{proof}
  \begin{enumerate}
  \item $\Gamma \Proves[\Sigma] !A$ पण $!A \notin
    \Gamma$ असे समजा. $\Gamma$ संपूर्ण
    $\Sigma$-सुसंगत असल्याने
    $\lnot!A \in \Gamma$. मग $\Gamma$ विसंगत
    ठरेल, कारण $!A, \lnot !A \Proves[\Sigma] \lfalse$.

  \item $!A \in \Sigma$ असेल, तर
    $\Gamma \Proves[\Sigma] !A$; आणि
    निगामी बंदतेवरून, म्हणजे
    \olref{prop:ccs-closed} वरून, $!A
    \in \Gamma$.

  \tagitem{prvFalse}{$\lfalse \in \Gamma$ असेल,
    तर $\Gamma \Proves[\Sigma] \lfalse$;
    त्यामुळे $\Gamma$ हा $\Sigma$-विसंगत
    ठरेल.}{}

  \tagitem{prvTrue}{$\Gamma \Proves[\Sigma] \ltrue$;
    म्हणून निगामी बंदतेवरून, म्हणजे
    \olref{prop:ccs-closed} वरून,
    $\ltrue \in \Gamma$.}{}

\item $\lnot!A \in \Gamma$ असेल, तर
    सुसंगततेमुळे $!A \notin \Gamma$.
    उलट $!A \notin \Gamma$ असेल, तर
    $\Gamma$ संपूर्ण $\Sigma$-सुसंगत
    असल्याने $\lnot!A \in \Gamma$.

  \tagitem{prvAnd}{\iftag{probAnd}{सराव.}{$!A \land !B \in
      \Gamma$ असे समजा. $(!A \land !B) \lif !A$
      हा सर्वतः-सत्य दृष्टांत असल्याने,
      निगामी बंदतेवरून, म्हणजे
      \olref{prop:ccs-closed} वरून, $!A \in \Gamma$.
      $!B \in \Gamma$ साठीही तसेच. उलट
      $!A \in \Gamma$ आणि $!B \in
      \Gamma$ दोन्ही असतील, तर निगामी
      बंदतेवरून $(!A \land !B) \in
      \Gamma$; कारण $!A \lif (!B \lif (!A \land !B))$
      हा सर्वतः-सत्य दृष्टांत आहे.}}{}

  \tagitem{prvOr}{\iftag{probOr}{सराव.}{$!A \lor !B \in
      \Gamma$, पण $!A \notin \Gamma$ आणि
      $!B \notin \Gamma$ असे समजा. $\Gamma$
      संपूर्ण $\Sigma$-सुसंगत असल्याने
      $\lnot!A \in \Gamma$ आणि
      $\lnot !B \in \Gamma$. मग
      $\lnot(!A \lor !B) \in \Gamma$;
      कारण $\lnot !A \lif (\lnot !B \lif \lnot (!A \lor !B))$
      हा सर्वतः-सत्य दृष्टांत आहे. त्यामुळे
      $\Gamma$ हा $\Sigma$-विसंगत ठरेल,
      हा व्याघात आहे.}}{}

  \tagitem{prvIf}{\iftag{probIf}{सराव.}{$!A \lif !B \in
      \Gamma$ आणि $!A \in \Gamma$ असे समजा.
      मग $\Gamma \Proves[\Sigma] !B$;
      त्यामुळे निगामी बंदतेवरून $!B \in \Gamma$.
      उलट $!A
      \lif !B \notin \Gamma$ असेल, तर
      $\Gamma$ संपूर्ण $\Sigma$-सुसंगत
      असल्याने $\lnot (!A \lif !B) \in \Gamma$.
      $\lnot(!A \lif !B) \lif !A$ हा
      सर्वतः-सत्य दृष्टांत असल्याने निगामी
      बंदतेवरून $!A \in
      \Gamma$. तसेच $\lnot(!A \lif !B) \lif
      \lnot !B$ हा सर्वतः-सत्य दृष्टांत असल्याने
      $\lnot !B \in
      \Gamma$. $\Gamma$ हा $\Sigma$-सुसंगत
      असल्यामुळे $!B \notin \Gamma$.}}{}

  \tagitem{prvIff}{\iftag{probIff}{सराव.}{$!A \liff !B \in
      \Gamma$ असे समजा. $!A \in \Gamma$
      असेल, तर $!B \in \Gamma$; कारण $(!A
      \liff !B) \lif (!A \lif !B)$ हा
      सर्वतः-सत्य दृष्टांत आहे. तसेच
      $!B \in \Gamma$ असेल, तर $!A \in \Gamma$.
      म्हणून $!A \in \Gamma$ आणि $!B \in \Gamma$
      दोन्ही असतील, किंवा $!A \in \Gamma$
      आणि $!B \in \Gamma$ यांपैकी कोणतेही नसेल.

      उलट $!A \liff !B \notin \Gamma$
      असे समजा. $\Gamma$ संपूर्ण
      $\Sigma$-सुसंगत असल्याने
      $\lnot (!A \liff !B) \in
      \Gamma$. $\lnot(!A \liff !B) \lif (!A \lif \lnot !B)$
      हा सर्वतः-सत्य दृष्टांत असल्याने
      $!A \in \Gamma$ असेल, तर
      $\lnot !B \in
      \Gamma$; आणि $\Gamma$ हा
      $\Sigma$-सुसंगत असल्याने $!B \notin
      \Gamma$. तसेच $!B \in \Gamma$
      असेल, तर $!A \notin \Gamma$.
      म्हणून $!A \in \Gamma$ आणि
      $!B \in \Gamma$ दोन्ही नसतील,
      तसेच $!A \notin
      \Gamma$ आणि $!B \notin \Gamma$
      दोन्हीही नसतील.}}{}
  \end{enumerate}
\end{proof}

\begin{probtag}{probAnd,probOr,probIf,probIff}
  \olref[nml][com][ccs]{prop:ccs-properties}
  ची सिद्धता पूर्ण करा.
\end{probtag}

\end{document}
